\begin{longtable}{|p{3.5cm}|p{5.5cm}|p{6cm}|}

\hline

Categoria & Métodos e estratégias de avaliação identificados na literatura & Exemplos de estudos \\ \hline

\endfirsthead

\hline

Categoria & Métodos e estratégias de avaliação identificados na literatura & Exemplos de estudos \\ \hline

\endhead

\hline

\endfoot

\endlastfoot

Testes de usabilidade com usuários &

Testes de usabilidade presenciais, testes comparativos, avaliação de desempenho baseada em tarefas, realização de tarefas específicas, testes com protótipos e avaliação de interfaces por usuários finais. &

Usability Assessment of Secure Messaging for Clinical Document Sharing between Health Care Providers and Patients; Usability Testing by Older Americans of a Prototype Google Map Web Site To Select Nursing Homes; Redesigning web sites for older adults. \\ \hline

Questionários e escalas de avaliação &

Questionários estruturados, questionários de satisfação, escalas Likert, System Usability Scale (SUS), PSSUQ, e-Qual, VisAWI e instrumentos específicos desenvolvidos para avaliação da interface ou da experiência do usuário. &

Usability Assessment of Secure Messaging for Clinical Document Sharing between Health Care Providers and Patients; Old and Young Users' White Space Preferences for Online News Web Pages; Novice User Perception of e-Services: A Study in the Egyptian Public Sector. \\ \hline

Observação e análise de tarefas &

Observação direta dos participantes, análise de tarefas, análise de desempenho, registro das ações realizadas durante as tarefas, análise de logs de interação e análise de sessões de uso. &

Usability Study to Promote Co-Creation among People with Disabilities, Developers, and Makers with a Focus on the Assistive Technology Open Platform in Korea; Addressing the challenge of visual information access from digital image and video libraries; CodeA11y: Making AI Coding Assistants Useful for Accessible Web Development. \\ \hline

Think-aloud e protocolos verbais &

Think-aloud, protocolos verbais durante a realização das tarefas e análise dos comentários dos participantes durante a interação com o sistema. &

Usability Assessment of Secure Messaging for Clinical Document Sharing between Health Care Providers and Patients; Usability Study to Promote Co-Creation among People with Disabilities, Developers, and Makers with a Focus on the Assistive Technology Open Platform in Korea; Investigating the roles of knowledge and cognitive abilities in older adult information seeking on the Web. \\ \hline

Entrevistas e grupos focais &

Entrevistas individuais, entrevistas semiestruturadas, entrevistas retrospectivas, grupos focais, discussões com participantes e coleta de percepções dos usuários após a interação. &

Breaking Barriers: The Design and Development of an Assistive Technology Web App for Older Latinos with Disabilities in Daily Activities; Requirements elicitation with and for older adults; A usability study of websites for older travelers. \\ \hline

Avaliação por especialistas &

Avaliação heurística, revisão humana, avaliação por especialistas, validação por especialistas e análise de interfaces com base em diretrizes ou critérios previamente definidos. &

Redesigning web sites for older adults; Research-derived web design guidelines for older people; Evaluating websites for older adults: adherence to 'senior-friendly' guidelines and end-user performance. \\ \hline

Inspeção e walkthrough &

Heuristic Evaluation, Cognitive Walkthrough, Recoverability Walkthrough, inspeção baseada em metáforas e identificação e categorização de problemas de usabilidade. &

Recoverability Walkthrough: An Alternative to Evaluate Digital Inclusion Interfaces; Usability inspection by metaphors of human thinking compared to heuristic evaluation; A conceptual tool for usability problem identification in website development. \\ \hline

Avaliação de acessibilidade automatizada &

Ferramentas automatizadas de avaliação de acessibilidade, WAVE, A-Checker, validadores de marcação, validadores CSS e ferramentas de verificação de conformidade com padrões de acessibilidade. &

A methodological approach for designing and developing web-based inventories of mobile Assistive Technology applications; Accessible Web Design for Older Adults: Challenges and Solutions; CodeA11y: Making AI Coding Assistants Useful for Accessible Web Development. \\ \hline

Eye tracking e análise de interação visual &

Eye tracking, análise de fixações, análise de sacadas e observação do comportamento visual durante a interação com as interfaces. &

Age-Related Differences in Eye Tracking and Usability Performance: Website Usability for Older Adults; Accessible Web Design for Older Adults: Challenges and Solutions. \\ \hline

Experimentos controlados e estudos comparativos &

Experimentos controlados, experimentos within-subjects, comparação entre grupos de usuários, comparação entre versões de interfaces, pré-teste e pós-teste e estudos comparativos de desempenho. &

Age related differences and the depth vs. breadth tradeoff in hierarchical online information systems; Old and Young Users' White Space Preferences for Online News Web Pages; The effects of age and the design of web-based training on computer task performance. \\ \hline

Análise estatística dos resultados &

ANOVA, teste t, teste de Scheffe, teste de Mann-Whitney, teste de qui-quadrado, correlação de Spearman e análise de confiabilidade por alfa de Cronbach. &

Old and Young Users' White Space Preferences for Online News Web Pages; Novice User Perception of e-Services: A Study in the Egyptian Public Sector; A study on the acceptance of website interaction aids by older adults. \\ \hline

Análise qualitativa e análise de conteúdo &

Análise de conteúdo, análise temática, análise de incidentes críticos, análise qualitativa das respostas dos participantes e categorização de problemas encontrados durante a avaliação. &

Breaking Barriers: The Design and Development of an Assistive Technology Web App for Older Latinos with Disabilities in Daily Activities; Adding voices to support web navigation among a low digital literacy group; Usability Study to Promote Co-Creation among People with Disabilities, Developers, and Makers with a Focus on the Assistive Technology Open Platform in Korea. \\ \hline

\caption{Métodos e estratégias de avaliação de usabilidade e acessibilidade identificados nos estudos que respondem diretamente à RQ3.}

\label{tab:metodosRQ3}

\end{longtable}